\chapter{المسائل والمخططات}\label{ch:g5:data}

آخر فصول السنة يشغّل كل شيء: مسائل متعددة الخطوات
بالنقود والكميات والمدد، والمخططات التي تعرض المعلومة
بنظرة واحدة. والفكرة النجمة هي «لكل واحد» --- إيجاد ثمن
شيء واحد --- وستكبر لتصير التناسب في \cref{ch:g6:propdata}.

\section{المسائل المتعددة الخطوات}

\begin{method}[ترويض مسألة طويلة]\label{met:g5:data:steps}
\begin{enumerate}
  \item اقرأ إلى النهاية؛ وقل ما السؤال الأخير؛
  \item وخطّط بالعكس: ماذا أحتاج أولًا لأجيب عنه؟
  \item وحُلّ الخطوات بالترتيب، كاتبًا عملية واحدة وجملة
        قصيرة لكل خطوة؛
  \item وتحقّق من الجواب الأخير بالحسّ السليم وبتقدير.
\end{enumerate}
\end{method}

\begin{example}\label{ex:g5:data:multistep}
«تطلب مدرسة $4$ علب من $25$ كتابًا ثمن الواحد $6$، وتدفع
$180$ للتوصيل فوق ذلك. فما الكلفة الكلية؟»
\begin{enumerate}
  \item الكتب: $4 \times 25 = 100$ كتاب.
  \item وثمن الكتب: $100 \times 6 = 600$.
  \item \omterm{def:g1:addition:def}{والمجموع}: $600 + 180 = 780$.
\end{enumerate}
والتحقّق بالتقدير: نحو مئة كتاب ثمن الواحد $6$ يقارب $600$، زائد
نحو $200$: أي نحو $800$ --- فالجواب $780$ معقول.
\end{example}

\begin{example}[لكل واحد]\label{ex:g5:data:perone}
«$5$ أكواب متماثلة ثمنها $12.50$؛ فكم يكلّف $8$ أكواب؟»
جِد ثمن كوب \emph{واحد} أولًا:
\[
12.50 \div 5 = 2.50,
\qquad\text{ثم}\qquad
8 \times 2.50 = 20 .
\]
فمن خلال الواحد يصير كل شيء في المتناول --- وهذا النمط ذو
الخطوتين يحلّ عائلة كاملة من المسائل (الوصفات، والسرعات،
والأثمان)، وهو بذرة \cref{ch:g6:propdata}.
\end{example}

\begin{example}[الباقي والميزانية]\label{ex:g5:data:budget}
عند زينة $30$. وتشتري كتابًا ثمنه $12.40$ وقلمين ثمن الواحد $2.30$.
فهل تستطيع أيضًا أن تشتري لعبة لوحية ثمنها $13$؟
\begin{enumerate}
  \item الأقلام: $2 \times 2.30 = 4.60$.
  \item والمصروف حتى الآن: $12.40 + 4.60 = 17$.
  \item \omterm{def:g3:division:remainder}{والباقي}: $30 - 17 = 13$ --- وهو يكفي اللعبة بالضبط،
        ولا يبقى شيء.
\end{enumerate}
\end{example}

\section{قراءة المخططات}

\begin{method}[قراءة أي مخطط]\label{met:g5:data:read}
\begin{enumerate}
  \item اقرأ العنوان: ما الذي يُعدّ؟
  \item واقرأ المحاور أو المفتاح: ماذا يمثّل التدريج الواحد، أو
        العمود الواحد، أو الرمز الواحد؟
  \item وعندئذٍ فقط أجب عن الأسئلة --- والإشارة إلى المخطط
        بمسطرة تساعد على قراءة الارتفاعات بدقة.
\end{enumerate}
\end{method}

\begin{example}\label{ex:g5:data:chart}
تساقط المطر في مدينة، شهرًا شهرًا:

\begin{omfigure}
\begin{tikzpicture}
\begin{axis}[
  width=8.8cm, height=5.0cm,
  ybar, bar width=13pt,
  symbolic x coords={Jan,Feb,Mar,Apr,May,Jun},
  xticklabels={يناير,فبراير,مارس,أبريل,مايو,يونيو},
  xtick=data,
  ymin=0, ymax=90,
  ytick={20,40,60,80},
  ylabel={المطر (مم)},
  tick label style={font=\small},
  label style={font=\small}]
  \addplot[fill=omDef!30, draw=omDef] coordinates
    {(Jan,60) (Feb,45) (Mar,55) (Apr,70) (May,80) (Jun,30)};
\end{axis}
\end{tikzpicture}

{\small مخطط بالأعمدة لتساقط المطر. التدريج الواحد $= 20$ مم.}
\end{omfigure}

القراءات: أمطر الشهور هو مايو ($80$ مم)؛ ونال يونيو نصف
مطر أبريل ($30$ مقابل $70$؟ لا --- فنصف $70$ هو $35$، إذن
\emph{ليس نصفه تمامًا}: فالقراءة الدقيقة تهمّ!)؛ \omterm{def:g1:addition:def}{والمجموع} في
الشهور الستة: $60 + 45 + 55 + 70 + 80 + 30 = 340$ مم.
\end{example}

\begin{example}[كمية متغيرة]\label{ex:g5:data:line}
ارتفاع نبتة، مقيسًا كل أسبوع:

\begin{omfigure}
\begin{tikzpicture}
\begin{axis}[omaxis,
  width=8.4cm, height=5.0cm,
  xmin=0, xmax=6.5, ymin=0, ymax=24,
  xlabel={الأسبوع}, ylabel={الارتفاع (سم)},
  xtick={1,...,6}, ytick={5,10,15,20}]
  \addplot[omDef, thick, mark=*, mark size=1.6pt] coordinates
    {(1,3) (2,6) (3,10) (4,15) (5,19) (6,21)};
\end{axis}
\end{tikzpicture}

{\small وصل نقاط القياس يُظهر \emph{النمو}: فكلما اشتدّ ميل
القطعة كان نمو النبتة في ذلك الأسبوع أسرع.}
\end{omfigure}

نمت النبتة أسرع ما نمت بين الأسبوعين $3$ و $4$ ($+5$ سم) وتباطأت
في النهاية ($+2$ سم في الأسبوع $6$).
\end{example}

\section{تمارين}

\begin{exercise}[$\star$]\label{exo:g5:data:1}
حُلّ بالخطوات: «يستأجر نادٍ حافلة بمبلغ $240$ ويشتري $32$ تذكرة
متحف ثمن الواحدة $7$. فما الكلفة الكلية للرحلة؟»
\end{exercise}

\begin{exercise}[$\star$]\label{exo:g5:data:2}
حُلّ بطريقة «لكل واحد»: «$3$ كغ من التفاح ثمنها $7.50$. فكم يكلّف
$5$ كغ؟»
\end{exercise}

\begin{exercise}[$\star$]\label{exo:g5:data:3}
حُلّ: «يتقاسم ستة أصدقاء بالتساوي كلفة نزهة قدرها $87$. فكم
يدفع كل واحد؟»
\end{exercise}

\begin{exercise}[$\star$]\label{exo:g5:data:4}
عند سامي $25$. ويشتري قميصًا ثمنه $13.90$ وقبعة ثمنها $8.50$. فكم
من المال بقي له؟
\end{exercise}

\begin{exercise}[$\star$]\label{exo:g5:data:5}
بالاستعانة بمخطط المطر في \cref{ex:g5:data:chart}: أي الشهور
نالت أكثر من $50$ مم؟ وبكم يزيد المطر في مايو على مطر
يونيو؟
\end{exercise}

\begin{exercise}[$\star$]\label{exo:g5:data:6}
بالاستعانة بمخطط النبتة في \cref{ex:g5:data:line}: كم كان ارتفاع
النبتة في الأسبوع $2$؟ وبين أي أسبوعين نمت
$4$ سم بالضبط؟
\end{exercise}

\begin{exercise}[$\star$]\label{exo:g5:data:7}
اُرسم مخططًا بالأعمدة لعدد الأهداف التي سجّلها فريق:
سبتمبر $4$، وأكتوبر $7$، ونوفمبر $3$، وديسمبر $6$. واختر
التدريج، وأعطِ \omterm{def:g1:addition:def}{المجموع}.
\end{exercise}

\begin{exercise}[$\star$]\label{exo:g5:data:8}
رزمة من $8$ علب لبن ثمنها $3.20$؛ وتُباع العلبة وحدها بثمن $0.50$.
احسب ثمن «الواحدة» في الرزمة، وكم يوفّر المرء في كل علبة
بشراء الرزمة.
\end{exercise}

\begin{exercise}[$\star\star$]\label{exo:g5:data:9}
«باعت سينما $148$ تذكرة ثمن الواحدة $9.50$ يوم السبت و $95$ تذكرة
بالثمن نفسه يوم الأحد. فبكم يزيد ما جنته يوم السبت على ما جنته
يوم الأحد؟» حُلّها بطريقتين: بحساب مبلغَي اليومين، أو
بالحساب من \emph{\omterm{ex:g1:subtraction:difference}{فرق}} التذاكر --- وتحقّق أن
الجوابين متفقان.
\end{exercise}

\begin{exercise}[$\star\star$]\label{exo:g5:data:10}
وصفة تكفي $4$ أشخاص تحتاج إلى $300$ غ من الأرز. وتطبخ عالية
لعدد $10$ أشخاص. فكم أرزًا تحتاج؟ (لكل شخص واحد أولًا ---
والأعداد العشرية مسموحة.)
\end{exercise}

\begin{exercise}[$\star\star$]\label{exo:g5:data:11}
اِبتكر مسألة من ثلاث خطوات تستعمل النقود جوابها الأخير
$5.50$، واكتب حلّها الكامل على أسلوب
\cref{met:g5:data:steps}، واختبرها على زميل.
\end{exercise}
